%=====================================================================
\chapter{Input Size, Input Nature, and the Three Cases}\label{ch:cases}
%=====================================================================
\locator{1}{Part I --- Analysing an Algorithm}

\begin{outcomes}
By the end of this chapter you will be able to:
\begin{tight}
  \item \textbf{Define} the size of an input for an array, a matrix, a graph
        and a number, and \textbf{justify} why a number's size is its bit
        count.
  \item \textbf{Explain} what is meant by the best, expected and worst case
        behaviour of an algorithm \emph{(CLO-1)}.
  \item \textbf{Identify} the characteristics of the data, or the other
        conditions and assumptions, that lead an algorithm to different
        behaviours \emph{(CLO-2)}.
  \item \textbf{Derive} the best, worst and average case running time of
        insertion sort and linear search, and \textbf{say} which to report and
        why.
  \item \textbf{Distinguish} an average taken over inputs from an expectation
        taken over an algorithm's own random choices.
\end{tight}
\end{outcomes}

\begin{prereq}
\chref{role} --- the RAM model and the idea of an order of growth. You will
also need to be comfortable that $\lg$ means $\log_{2}$, and that
$\sqrt{2^{b}} = 2^{b/2}$.
\end{prereq}

\begin{keyterms}
input size $\bullet$ bit length $\bullet$ sparse $\bullet$ dense $\bullet$
best case $\bullet$ worst case $\bullet$ average case $\bullet$ expected
running time $\bullet$ input distribution $\bullet$ uniform random permutation
$\bullet$ randomised algorithm $\bullet$ adversary $\bullet$
pseudo-polynomial
\end{keyterms}

\section{The Size of the Input}

Everything in \chref{role} was written as a function of $n$ without ever saying
what $n$ counts. For an array that is obvious enough. For the other things
algorithms take as input it is a decision, and a decision that changes the
answer.

\begin{definitionbox}[Input Size]
The \term{input size} is whatever parameter the running time is most naturally
expressed in --- but it must be a measure of \emph{how much input there is},
not of what the input says.

\medskip
\rowcolors{2}{white}{lightbg}
\begin{longtable}{@{}L{3.0cm}L{11.5cm}@{}}
\toprule
\rowcolor{primary!15}
\textbf{Input} & \textbf{Size, by convention} \\
\midrule
\endhead
\textbf{array, list} & $n$, the number of elements \\
\textbf{string} & $n$, the number of characters \\
\textbf{$n \times n$ matrix} & $n$, the side --- so an $\bigO{n^{3}}$ algorithm
  is \emph{linear} in the $n^{2}$ numbers it reads, which is worth saying out
  loud \\
\textbf{graph} & \emph{two} parameters: $|V|$ vertices and $|E|$ edges \\
\textbf{number} & $b$, the number of \emph{bits} needed to write it \\
\bottomrule
\end{longtable}
\end{definitionbox}

Two of those rows need an argument.

\subsection*{Graphs take two parameters, and that is not a technicality}

A graph has $|V|$ vertices and $|E|$ edges, and $|E|$ can be anywhere from $0$
to $|V|(|V|-1)/2$. A bound of $\bigO{|V| + |E|}$ and a bound of
$\bigO{|V|^{2}}$ are the same thing on a \term{dense} graph and wildly
different on a \term{sparse} one.

\begin{worked}[The Same Bound, Two Graphs]
Breadth-first search is $\bigO{|V| + |E|}$ from an adjacency list.

\smallskip
\textbf{A road network:} 2 million junctions, and each junction meets about
three roads, so $|E| \approx 3 \times 10^{6}$.
\[ |V| + |E| \approx 2\!\times\!10^{6} + 3\!\times\!10^{6} = 5 \times 10^{6}. \]

\smallskip
\textbf{A social graph of the same vertex count}, each person connected to
500 others, so $|E| \approx 2\!\times\!10^{6} \times 500 / 2 = 5 \times
10^{8}$.
\[ |V| + |E| \approx 5 \times 10^{8}. \]

\smallskip
\textbf{A factor of 100}, on the same algorithm, at the same $|V|$. This is
why \chref{graphs} comes before the graph algorithms rather than after them,
and why every bound in Part~V is quoted in both parameters.
\end{worked}

\subsection*{A number's size is its length, not its value}

This one students get wrong, and it is the root of a famous confusion that
\chref{dptwo} and \chref{complexity} both depend on.

\begin{alertbox}[Why not just use $N$ itself?]
Because the input to a program is \emph{written down}, and a number $N$ is
written in about $b = \lfloor \lg N \rfloor + 1$ bits. Typing one more digit
multiplies $N$ by ten while adding three and a bit bits of input.

\smallskip
If we measured size by value, then an algorithm that loops $N$ times would be
``linear'' --- and we would be calling an algorithm efficient when doubling the
\emph{length} of its input multiplies its running time by $2^{b}$.

\smallskip
So: \textbf{size $= b \approx \lg N$}, and an algorithm that loops $N$ times is
$\bigO{2^{b}}$ --- exponential. An algorithm whose cost is polynomial in $N$ but
exponential in $b$ is called \term{pseudo-polynomial}; the 0/1 knapsack
algorithm of \chref{dptwo} is the standard example, and \chref{complexity}
explains why that distinction decides whether a problem counts as solved.
\end{alertbox}

\begin{worked}[The Claim, Measured]
\texttt{code/ch02\_cases.cpp} tests a prime by trial division --- a loop running
about $\sqrt{N}/2$ times. Since $\sqrt{N} = \sqrt{2^{b}} = 2^{b/2}$, adding
\emph{two} bits should exactly double the work. Each row below is a prime two
bits longer than the row above; median of three runs.

\rowcolors{2}{white}{lightbg}
\begin{center}
\begin{tabular}{@{}rrrr@{}}
\toprule
\rowcolor{primary!15}
bits & $N$ & time (ms) & ratio \\
\midrule
32 & 4{,}294{,}967{,}291     &  0.22 & ---  \\
34 & 17{,}179{,}869{,}143    &  0.44 & 2.01 \\
36 & 68{,}719{,}476{,}731    &  0.86 & 1.94 \\
38 & 274{,}877{,}906{,}899   &  1.76 & 2.06 \\
40 & 1{,}099{,}511{,}627{,}689 & 3.63 & 2.07 \\
42 & 4{,}398{,}046{,}511{,}093 & 6.93 & 1.91 \\
44 & 17{,}592{,}186{,}044{,}399 & 15.75 & 2.27 \\
46 & 70{,}368{,}744{,}177{,}643 & 29.83 & 1.89 \\
48 & 281{,}474{,}976{,}710{,}597 & 58.81 & 1.97 \\
\bottomrule
\end{tabular}
\end{center}

\smallskip
\textbf{Eight ratios, mean 2.02.} The prediction is exact and the measurement
meets it.

\smallskip
\textbf{Now extrapolate, which is the point.} From 32 bits to 48 the time went
from 0.22 ms to 58.81 ms --- a factor of 267, against a predicted $2^{8} =
256$. Carry that on:
\begin{itemize}[leftmargin=1.6em, itemsep=1pt, topsep=2pt]
  \item \textbf{64 bits:} $2^{8} \times 58.81\,\text{ms} \approx 15$ seconds.
  \item \textbf{128 bits:} $2^{40} \times 58.81\,\text{ms} \approx 6.5 \times
        10^{10}$ s $\approx$ \textbf{2{,}000 years}.
  \item \textbf{2048 bits}, which is what an RSA key is: there is no useful way
        to write the number.
\end{itemize}

\smallskip
\textbf{And notice what did \emph{not} happen.} Going from 32 bits to 128 bits
is going from a four-byte input to a sixteen-byte input --- twelve extra bytes.
The input grew by a factor of four and the running time by a factor of $4
\times 10^{11}$. That is what ``exponential in the input size'' means, and it
is why measuring the size by the value would have hidden it completely.
\end{worked}

\section{The Nature of the Input: Three Cases}

Fix the size. Two inputs of the same size can still cost an algorithm wildly
different amounts, because an algorithm's control flow depends on what the data
\emph{says}, not only on how much of it there is.

\begin{definitionbox}[Best, Worst and Average Case]
Let $T(x)$ be the number of primitive operations the algorithm performs on
input $x$, and let $X_{n}$ be the set of all inputs of size $n$.

\medskip
\rowcolors{2}{white}{lightbg}
\begin{longtable}{@{}L{4.6cm}L{9.9cm}@{}}
\toprule
\rowcolor{primary!15}
\textbf{Case} & \textbf{What it promises} \\
\midrule
\endhead
\textbf{best:}\quad $\displaystyle\min_{x \in X_{n}} T(x)$
  & Nothing useful --- a lower bound on an optimist's hopes. \\
\textbf{worst:}\quad $\displaystyle\max_{x \in X_{n}} T(x)$
  & A \emph{guarantee}: no input of size $n$ costs more than this. \\
\textbf{average:}\quad $\displaystyle\mathop{\mathbb{E}}_{x \sim D}
  \bigl[T(x)\bigr]$
  & What to expect \emph{if} inputs really are drawn from $D$ --- and $D$ must
    be stated. \\
\bottomrule
\end{longtable}
\end{definitionbox}

\begin{alertbox}[``Average case'' is meaningless until you name a distribution]
There is no such thing as a random input. There is only an input drawn from a
distribution you have chosen, and the choice is a modelling assumption that can
be wrong.

\smallskip
The conventional assumption for sorting is a \term{uniform random
permutation}: every ordering of the $n$ elements equally likely. It is a
reasonable default and it is often false --- real data arrives partly sorted,
or sorted by the wrong key, or in a few long runs. When you quote an average
case, say what distribution it averages over. When you read one, check.
\end{alertbox}

\subsection*{Insertion sort, in all three cases}

Insertion sort's inner loop shifts elements right until it finds the place for
\texttt{key}. The number of shifts for element $j$ is how many of the $j$
elements before it are larger --- the number of \emph{inversions} involving
$j$.

\begin{worked}[The Three Cases, Derived]
\textbf{Best case --- already sorted.} For every $j$, $a[j-1] \le \texttt{key}$,
so the \texttt{while} test fails at once and does zero shifts. The outer loop
still runs $n-1$ times:
\[ T(n) = c(n-1) = \bigTh{n}. \]

\smallskip
\textbf{Worst case --- strictly decreasing.} Every element must travel to the
front, so element $j$ does $j$ shifts:
\[ \sum_{j=1}^{n-1} j = \frac{(n-1)n}{2} \approx \frac{n^{2}}{2}
   = \bigTh{n^{2}}. \]

\smallskip
\textbf{Average case --- uniform random permutation.} For any pair $(i,j)$ with
$i<j$, the two orderings are equally likely, so the pair is inverted with
probability $\tfrac12$. There are $\binom{n}{2}$ pairs, so by linearity of
expectation the expected number of inversions is
\[ \mathbb{E}[\text{inversions}] = \frac12\binom{n}{2}
   = \frac{n(n-1)}{4} \approx \frac{n^{2}}{4} = \bigTh{n^{2}}. \]

\smallskip
\textbf{The conclusion, and it is the one that matters.} The average case is
\emph{exactly half} the worst case. A factor of two is a constant, and
$\bigTh{\cdot}$ discards constants --- so insertion sort's average and worst
cases have the \emph{same order of growth}, and the only case that differs is
the best one.
\end{worked}

\begin{worked}[One Algorithm, Three Orders of Growth, Measured]
\texttt{code/ch02\_cases.cpp} runs the same insertion sort on three shapes of
input at each size; median of five runs, times in milliseconds.

\rowcolors{2}{white}{lightbg}
\begin{center}
\begin{tabular}{@{}rrrrrrr@{}}
\toprule
\rowcolor{primary!15}
$n$ & sorted & ratio & random & ratio & reversed & ratio \\
\midrule
 2{,}000 & 0.002 & ---  &  0.340 & ---  &   0.634 & ---  \\
 4{,}000 & 0.004 & 1.90 &  1.430 & 4.21 &   2.551 & 4.03 \\
 8{,}000 & 0.008 & 2.05 &  5.288 & 3.70 &   9.736 & 3.82 \\
16{,}000 & 0.021 & 2.73 & 21.357 & 4.04 &  42.714 & 4.39 \\
32{,}000 & 0.028 & 1.33 & 89.790 & 4.20 & 185.054 & 4.33 \\
\bottomrule
\end{tabular}
\end{center}

\smallskip
\textbf{The sorted column doubles} --- 1.90, 2.05, 2.73, 1.33. The ratios are
ragged because the absolute times are two to twenty-eight \emph{microseconds},
where the clock's own resolution and one stray interrupt are a large fraction
of the measurement. The honest reading is: these are five numbers consistent
with linear growth and far too small to be consistent with quadratic, which is
all the experiment is being asked to show.

\smallskip
\textbf{The random and reversed columns quadruple}, mean 4.04 and 4.14, as
$\bigTh{n^{2}}$ requires.

\smallskip
\textbf{And now the prediction the derivation made.} Random should be half of
reversed, because $n^{2}/4$ is half of $n^{2}/2$. Measured:
\[ 0.54,\quad 0.56,\quad 0.54,\quad 0.50,\quad 0.49. \]
Five independent confirmations of a factor the theory produced from counting
inversions. This is what it looks like when an analysis is right about more
than the exponent.

\smallskip
\textbf{Finally, the spread.} At $n = 32{,}000$: 0.028 ms on sorted input,
185 ms on reversed. Same algorithm, same compiler, same machine, same input
size --- a factor of \textbf{6{,}600}. ``How big is the input'' does not come
close to determining how long a program takes.
\end{worked}

\section{Which Case Should You Report?}

\begin{conceptbox}[The convention, and the reasons for it]
\textbf{Report the worst case unless you say otherwise.} Three reasons:

\begin{tightnum}
  \item It is a \emph{guarantee}. An average tells you about a population of
        runs; a worst case tells you about the run you are having now, which is
        the one with the deadline.
  \item It is often what happens. For searching an absent key, or any
        algorithm that must examine everything to be sure, the worst case is
        the typical case.
  \item It is frequently the same order as the average anyway --- insertion
        sort above, where the factor between them is two.
\end{tightnum}

\smallskip
\textbf{Report the best case almost never.} ``This sort is $\bigO{n}$'' with the
best case quietly intended is the most common dishonest sentence in the
subject.

\smallskip
\textbf{Report the average case when the worst case is both rare and
pessimistic}, as with quicksort (\chref{quicksort}) and hashing
(\chref{hashing}) --- and then state the distribution, because in both of those
chapters there is an input that makes the worst case happen on purpose.
\end{conceptbox}

\begin{definitionbox}[Expected Running Time]
A \term{randomised algorithm} makes random choices of its own. Its
\term{expected running time} on an input $x$ is the average over
\emph{those choices}, with $x$ held fixed:
\[ \mathbb{E}\bigl[T(x)\bigr], \qquad
   \text{expectation over the algorithm's coin flips.} \]
\end{definitionbox}

\begin{alertbox}[Two different averages, and the difference is the whole trick]
\textbf{Average-case analysis} averages over inputs. You must assume a
distribution on the input, and an adversary who knows your algorithm can choose
an input outside it.

\smallskip
\textbf{Expected-time analysis} of a randomised algorithm averages over the
algorithm's own coins, \emph{for every input}. There is no assumption about the
data at all. An adversary who knows your code still cannot pick a bad input,
because he does not know your coins.

\smallskip
That is why \chref{quicksort} randomises: it converts an assumption about the
world, which can be violated, into an assumption about a random number
generator, which cannot. The bound goes from ``$\bigTh{n \lg n}$ \emph{if} your
data is unsorted'' to ``$\bigTh{n \lg n}$ expected, on \emph{every} input''.
\end{alertbox}

\begin{pitfall}
\begin{tight}
  \item \textbf{Quoting one case and meaning another.} Always name it: ``linear
        search is $\bigTh{n}$ in the worst case and $\bigTh{n}$ on average'' is
        a real statement; ``linear search is $\bigO{n}$'' is a weaker one that
        people read as the worst case, correctly.
  \item \textbf{Measuring the size of a number by its value.} Then trial
        division looks linear, and the table in Section~2.1 looks impossible.
  \item \textbf{Quoting a graph bound in one parameter.} $\bigO{|V|^{2}}$ hides
        whether you have exploited sparsity; $\bigO{|V| + |E|}$ says you have.
  \item \textbf{Assuming the average case without saying so.} Especially for
        sorting, where real inputs are very often partly sorted and the uniform
        permutation assumption is simply false.
  \item \textbf{Thinking the best case tells you the algorithm is fast.}
        Insertion sort's best case is $\bigTh{n}$ and it is still a quadratic
        sort.
  \item \textbf{Confusing ``average case'' with ``expected time''.} The first
        assumes something about the data; the second assumes nothing. The
        second is much stronger and is almost always what a randomised
        algorithm's bound means.
\end{tight}
\end{pitfall}

\begin{lab}[Three Shapes, and the Cost of a Bit]
Reproduces both tables in this chapter. Expect twenty-five minutes --- part B
takes about two minutes of CPU on its own.

\begin{lstlisting}[language=C++]
// ch02_cases.cpp -- the nature of the input, and the size of the input.
// Build: cl /O2 /EHsc /std:c++17 ch02_cases.cpp      (see Appendix A)

// --- 1. Three shapes of input, all of the same size -----------------
// Best case: already sorted, so the inner while never executes.
// Worst case: reversed, so it runs to the front every time.
// Average: a random permutation, which moves each element about half
// way -- and half of a quadratic is still a quadratic.
static std::vector<int> sorted_input(size_t n) {
    std::vector<int> a(n);
    std::iota(a.begin(), a.end(), 0);
    return a;
}
static std::vector<int> reversed_input(size_t n) {
    std::vector<int> a = sorted_input(n);
    std::reverse(a.begin(), a.end());
    return a;
}

// --- 2. Trial division, whose cost is governed by sqrt(N) -----------
// The loop runs about sqrt(N)/2 times when n is prime, which is the
// worst case for this algorithm -- so a prime is what we time.
static bool is_prime_trial(uint64_t n) {
    if (n < 2) return false;
    if (n % 2 == 0) return n == 2;
    for (uint64_t d = 3; d * d <= n; d += 2)
        if (n % d == 0) return false;
    return true;
}

// --- 3. The same algorithm, three orders of growth ------------------
for (size_t n = 2000; n <= 32000; n *= 2) {
    auto s = sorted_input(n);
    auto r = random_input(n, 12345);
    auto v = reversed_input(n);
    double ts = time_ms([&]{ auto a = s; insertion_sort(a); }, 5);
    double tr = time_ms([&]{ auto a = r; insertion_sort(a); }, 5);
    double tv = time_ms([&]{ auto a = v; insertion_sort(a); }, 5);
    /* print the three times and the three ratios */
}

// --- 4. Why the size of a number is its number of bits --------------
// Each row adds TWO bits -- four more characters of input, if you
// measure input by its value. Watch the time double anyway.
const uint64_t primes[] = {
    4294967291ULL,      17179869143ULL,     68719476731ULL,
    274877906899ULL,    1099511627689ULL,   4398046511093ULL,
    17592186044399ULL,  70368744177643ULL,  281474976710597ULL};
int bits = 32;
for (uint64_t p : primes) {
    bool r = false;
    double t = time_ms([&]{ r = is_prime_trial(p); }, 3);
    if (!r) { printf("  BUG: %llu is not prime\n", p); return 1; }
    printf("%4d %22llu %12.2f %7.2f\n", bits, p, t, t / prev);
    prev = t; bits += 2;
}
\end{lstlisting}

\textbf{What to notice.} In part A, look at the \emph{ratio} of the random
column to the reversed column rather than at either alone. It should be close to
$0.5$, and it is the one number in the experiment that tests the inversion
count rather than just the exponent.

\smallskip
\textbf{Then try to break it.} Replace the random permutation with an array
that is sorted except for the last 100 elements, which are shuffled. Predict
the running time before you measure it. This is the input shape real data
actually has, and it is the one both the best-case and the average-case
analyses describe badly.
\end{lab}

\begin{checkpoint}
\begin{tightnum}
  \item An algorithm takes as input an $n \times n$ matrix and runs in
        $\bigTh{n^{3}}$. Is it polynomial in the size of its input? Answer in
        terms of the number of numbers actually read.
  \item Linear search for a key that is present at a uniformly random position
        examines $(n+1)/2$ elements on average and $n$ in the worst case. Why
        does the course nevertheless describe both cases with the same
        $\bigTh{\cdot}$?
  \item A colleague reports that their sorting routine is ``$\bigO{n \lg n}$
        on average''. What single question would you ask before believing the
        number means anything for your data?
\end{tightnum}
\end{checkpoint}

\begin{chaptersummary}
\begin{tight}
  \item Input size is a measure of how much input there is: elements for an
        array, the side for a matrix, \emph{two} parameters for a graph, and
        \textbf{bits} for a number.
  \item A number's size is its bit length, so a loop running $N$ times is
        $\bigO{2^{b}}$ --- exponential. Measured: trial division doubled for
        every two bits added, eight times running, mean ratio 2.02.
  \item Best case promises nothing, worst case promises a guarantee, average
        case promises something only once you have named the distribution.
  \item Insertion sort: $\bigTh{n}$ best, $\bigTh{n^{2}}$ worst, $\bigTh{n^{2}}$
        average --- the average being exactly half the worst, which measurement
        confirmed at 0.54, 0.56, 0.54, 0.50, 0.49.
  \item At $n = 32{,}000$ the same sort took 0.028 ms or 185 ms depending only
        on the \emph{shape} of its input: a factor of 6{,}600.
  \item Report the worst case by default. Report an average only with its
        distribution attached.
  \item An average over inputs assumes something about the world; the expected
        time of a randomised algorithm assumes nothing, which is why
        \chref{quicksort} randomises.
\end{tight}
\end{chaptersummary}

\begin{reviewq}
  \item \textbf{(MCQ)} The input size of the integer $N$ is taken to be:
        (a)~$N$; (b)~$\sqrt{N}$; (c)~about $\lg N$; (d)~$N^{2}$.
  \item \textbf{(MCQ)} Insertion sort's average-case running time on a uniform
        random permutation is: (a)~$\bigTh{n}$; (b)~$\bigTh{n \lg n}$;
        (c)~$\bigTh{n^{2}}$; (d)~$\bigTh{n^{2}\lg n}$.
  \item \textbf{(Short)} Explain what is meant by the best, expected and worst
        case behaviour of an algorithm, and give one algorithm whose three
        cases have three different orders of growth \emph{(CLO-1)}.
  \item \textbf{(Short)} Identify two characteristics of input data that change
        which case an algorithm exhibits, and name an algorithm affected by
        each \emph{(CLO-2)}.
  \item \textbf{(Short)} Why is a bound of $\bigO{|V|+|E|}$ more informative
        than one of $\bigO{|V|^{2}}$, even when the two are equal for the graph
        in front of you?
  \item \textbf{(Applied)} Trial division on a 48-bit prime took 58.81 ms on
        the book's machine. Estimate the time for a 70-bit prime, state the
        assumption your estimate rests on, and say why the estimate is more
        trustworthy than the absolute time it is derived from.
  \item \textbf{(Applied)} You are given a sort you cannot read the source of.
        Design an experiment that decides whether its worst case is quadratic,
        using only the ability to supply inputs and time it. Say what inputs
        you would supply, what you would compute, and what result would leave
        you unable to tell.
\end{reviewq}
